\section{논의}
\label{sec:discussion}

\begin{revision}

\subsection{검사 가능성과 도구의 역할}
조건별 의무와 해제 증거를 연결하면, 새로운 근거 B가 앞선 A를 해제하는지와 A가 미해제인지 미상인지를 구분할 수 있다. 최초 문제 사건ㆍ관련 의무ㆍ사유를 원문 구절에 연결하는 출력은 수정할 주석이나 추가 증거가 필요한 항목을 검토하는 데 쓸 수 있다. 그러나 작성 프레임의 정확성을 모델 검사가 보장하지는 않는다. 잘못 연결한 조건 ID나 누락된 의무는 별도 의미 검토가 필요하다.

이전 이력 창과 전체 이력의 비교는 기억 길이가 일부 판정에 영향을 주었음을 보여 준다. FSM의 144/144와 보완 모델의 동일 사례 144/144는 이 범위의 논리를 단순 상태 기계로도 실행할 수 있다는 근거다. UPPAAL의 역할은 표현한 상태ㆍ조건 속성을 명시적으로 질의하고 진단 경로와 모델 파일을 보존하는 데 있다. 정확도 우위나 UPPAAL만의 필요성을 입증하지 않는다.

\subsection{주석 검토에서 구분할 조치}
확정 모순은 미해제가 확인된 의무와 행동을 함께 검토한다. 근거 부족은 A의 해제 여부를 B로 대신 채우지 않고 추가 주석/관측을 요청할 대상으로 둔다. 양립 판정도 명시된 의무에 대해서만 유효하다. P4는 주석을 수정할 근거가 아니라 검사 상태 처리의 고장을 점검할 근거다. 따라서 모든 경고를 자연 오류라고 하거나 자료를 자동 삭제하는 절차로 해석하지 않는다.

자연 자료의 낮은 일치도와 합의 사유 누락은 현재 기준을 실제 오류 정답으로 쓰기 어렵게 한다. 이번 작성 시나리오는 이 문제를 해결한 사람이 평가한 gold가 아니다. 자동 추출의 전 사례 미상 결과도 별도 한계로 남는다. 독립 자료와 실제 전문가 검토는 후속 검증이며, 현재의 적용 가능성 실증을 그런 성능으로 확대하지 않는다.
\end{revision}
